\section{Related work}
\label{sec:related}

\paragraph{Embedded continual learning.} CL methods are classically split into regularization-, replay-
and architecture-based approaches~\cite{DeLange2021Survey}. EWC~\cite{Kirkpatrick2017EWC} belongs to the
first family: a quadratic penalty weighted by the Fisher information protects the parameters that matter
for past tasks, without storing any data --- a valuable property under memory constraints. On MCUs,
\gls{TinyOL}~\cite{Ren2021TinyOL} demonstrated the feasibility of online learning of an additional layer, and
\gls{HDC}~\cite{Benatti2019HDC} that of very low-power non-neural learning.
LifeLearner~\cite{Kwon2023LifeLearner} and quantized latent-replay
platforms~\cite{Ravaglia2021QLRCL} address the memory/accuracy trade-off on embedded targets. Recent
\gls{TinyML} surveys~\cite{Capogrosso2023TinyML} point out that \emph{on-device} training remains
under-explored compared with inference alone, and work applied to predictive
maintenance~\cite{Hurtado2023CLPdM} is still mostly evaluated off-target.

\paragraph{Quantizing: three degrees of freedom, rarely crossed.} Quantization is often reduced to a
choice of \emph{format} --- FP32 to INT8, or Q15 fixed point. Two further degrees govern it just as much.
The \emph{moment} first~\cite{Jacob2018,Krishnamoorthi2018}: quantizing after training, in PTQ, is immediate but sensitive to scale calibration,
whereas quantizing during training, in QAT, costs a retraining and assumes access to the learning pipeline,
which is not always available in an industrial context. The \emph{depth} next, together with the
granularity that comes with it: going below 8 bits pays off only if the scale is estimated finely enough,
per channel rather than per tensor, and if storage actually follows the advertised format. These three
axes are most often studied separately; we vary them on \emph{one head, one dataset and one board}, which
makes them comparable rather than merely juxtaposed.

\paragraph{What is rarely verified: the three promises, separately.} Quantization is usually justified by
a threefold argument --- less memory, less time, less energy --- inherited from targets equipped with
dedicated integer units. On a general-purpose core with an \gls{FPU}, nothing guarantees that the three terms
behave alike, and reporting them jointly requires instrumentation that few studies put in place: measuring
the \emph{total} system memory rather than the weight footprint alone, timing the kernel item by item, and
integrating a real current at the probe. Our empirical contribution is to treat these three promises as
three distinct questions, measured on the same board: we show that the first holds but bears on a minor
item, and that the other two do not. Along the way we establish the collapse of naive PTQ of a decision
head, isolate its cause rather than masking it with upstream QAT, and demonstrate its recovery through a
calibrated kernel.
